\chapter{深度学习计算}

层把输入映成输出并携带参数；块把层或块再复合。
训练需要访问、初始化与共享这些参数，必要时延后到首次前向才确定形状。
自定义层给出新的映射，读写把参数与张量固化到存储，计算设备决定张量所在的空间。

\section{层和块}

单输出线性模型是映射 $f:\mathbb{R}^{d}\to\mathbb{R}$，带参数 $\bm{\theta}$。
一层接受向量、产生向量、并由可训练参数描述；整个多层感知机仍是同一类对象。
块（block）是可嵌套的计算单元：输入进前向映射，输出可改形状，梯度由反向传播给出。

\subsection{自定义块}

块必须实现前向 $f(\bm{x})$，并持有参数以便求导与初始化。
以单隐层感知机为例，设 $\bm{x}\in\mathbb{R}^{20}$，
\begin{align*}
\bm{h}
&=
\operatorname{ReLU}(\bm{W}_1\bm{x}+\bm{b}_1),
\qquad
\bm{W}_1\in\mathbb{R}^{256\times 20},\\
\bm{o}
&=
\bm{W}_2\bm{h}+\bm{b}_2,
\qquad
\bm{W}_2\in\mathbb{R}^{10\times 256}.
\end{align*}
即 $f=\ell_2\circ\operatorname{ReLU}\circ\ell_1$。
形状允许 $f(\bm{x})\in\mathbb{R}^{10}$ 与 $\bm{x}\in\mathbb{R}^{20}$ 不同。
反向传播给出 $\nabla_{\bm{W}_1}L$、$\nabla_{\bm{W}_2}L$ 及对输入的 $\nabla_{\bm{x}}L$。

\subsection{顺序块}

顺序复合定义为
\[
\bigl(\operatorname{Seq}(f_1,\ldots,f_n)\bigr)(\bm{x})
=
f_n\circ f_{n-1}\circ\cdots\circ f_1(\bm{x}).
\]
追加操作把 $f_{n+1}$ 接到链条末端。
参数集是各子块参数的不交并
\[
\operatorname{params}\bigl(\operatorname{Seq}(f_1,\ldots,f_n)\bigr)
=
\bigsqcup_{i=1}^{n}\operatorname{params}(f_i),
\]
这要求每个 $f_i$ 登记为子模块，而不能只放入普通列表。

\subsection{在前向传播函数中执行代码}

并非所有架构都是 $f_n\circ\cdots\circ f_1$。
前向中可插入控制流或与参数无关的常数。
设 $c$ 为常数参数（constant parameter），层
\[
f(\bm{x},\bm{w})=c\cdot\bm{w}^{\mathsf T}\bm{x}
\]
在 $\bm{w}$ 上可训练，$c$ 不进入 $\operatorname{params}(f)$。
亦允许按范数分支，例如
\[
\bm{h}
\leftarrow
\begin{cases}
\bm{h}/\|\bm{h}\|_2, & \|\bm{h}\|_2>1,\\
\bm{h}, & \text{否则,}
\end{cases}
\]
再进入后续线性层。只要这些运算可微或次可微，反向传播仍给出 $\nabla_{\bm{\theta}}L$。

\subsection{效率}

前向若含大量宿主解释与字典查找，加速器上的核必须等待。
总时间可写成
\[
T=T_{\mathrm{host}}+T_{\mathrm{device}}+T_{\mathrm{copy}}.
\]
当 $T_{\mathrm{host}}$ 因全局解释器锁而串行时，$T_{\mathrm{device}}$ 被阻塞。
把可向量化的映射尽量放进一次设备核，可减小 $T_{\mathrm{host}}/T$。

\subsection{小结}

块是可嵌套映射 $f$，顺序块实现 $f_n\circ\cdots\circ f_1$，并负责参数登记与反向传播。

\subsection{练习}

核验列表存储是否仍使 $\operatorname{params}$ 递归可收集，以及并联
$[f_1(\bm{x}),f_2(\bm{x})]$ 与重复实例 $f\circ\cdots\circ f$ 的形状。

\section{参数管理}

架构与超参数固定后，训练求解
\[
\bm{\theta}^{\ast}\in\operatorname*{argmin}_{\bm{\theta}}L(\bm{\theta}).
\]
随后需要读出 $\bm{\theta}$ 以便复用、保存或检查。
下面把访问、初始化与共享写成对 $\bm{\theta}$ 的运算。

\subsection{参数访问}

顺序模型中第 $\ell$ 层参数是命名张量，通常含权重 $\bm{W}^{(\ell)}$ 与偏置 $\bm{b}^{(\ell)}$。
名称在整网中唯一，即使层数达数百。

\subsubsection{目标参数}

参数对象同时携带值、梯度与状态。
未反向传播时 $\nabla_{\bm{\theta}}L$ 处于初始（通常为零）状态；
一次 $L.\mathrm{backward}$ 之后
\[
\frac{\partial L}{\partial \bm{W}^{(\ell)}},\qquad
\frac{\partial L}{\partial \bm{b}^{(\ell)}}
\]
才有定义。对参数的数值运算必须先取出底层张量。

\subsubsection{一次性访问所有参数}

嵌套块的参数树为
\[
\operatorname{params}(B)
=
\operatorname{local}(B)
\cup
\bigcup_{B'\in\operatorname{children}(B)}\operatorname{params}(B').
\]
一次性遍历给出全部 $(\text{name},\bm{\theta}_{\text{name}})$，而不必按层手工索引。

\subsubsection{从嵌套块收集参数}

若 $B=\operatorname{Seq}(B_1,B_2,\ldots)$ 且每个 $B_i$ 自身又是顺序块，
则名称按层级拼接。
访问“第一主块、第二子块、第一层”的偏置，对应路径索引下的 $\bm{b}$。
分层嵌套与嵌套列表索引同构。

\subsection{参数初始化}

默认把 $\bm{W}$、$\bm{b}$ 抽自由输入输出维决定的均匀区间。
自定义规则是从另一分布 $P$ 抽样
\[
\bm{W}\sim P.
\]

\subsubsection{内置初始化}

高斯初始化取
\[
w_{ij}\sim\mathcal{N}(0,0.01^{2}),\qquad b_i=0.
\]
亦可令全部元素为常数，例如 $\bm{W}=\bm{1}$。
不同块可使用不同规则：第一层 Xavier，第三层常数 $42$。
Xavier 对 $\bm{W}\in\mathbb{R}^{n_{\mathrm{out}}\times n_{\mathrm{in}}}$ 满足
\[
\operatorname{Var}(w)
=
\frac{2}{n_{\mathrm{in}}+n_{\mathrm{out}}}
\quad\text{（或均匀区间按该方差标定）}.
\]

\subsubsection{自定义初始化}

按元素独立抽样
\[
w
\sim
\begin{cases}
U(5,10), & \text{概率 }1/4,\\
0, & \text{概率 }1/2,\\
U(-10,-5), & \text{概率 }1/4.
\end{cases}
\]
该规则仍是合法的 $P$，只是把质量放在远离零的两侧与原点。

\subsection{参数绑定}

共享指定 $\bm{W}^{(3)}=\bm{W}^{(5)}=\bm{W}$ 为同一张量，而非值相等的两份副本。
于是 $\mathrm{d}\bm{W}^{(3)}=\mathrm{d}\bm{W}^{(5)}$。
反向传播时梯度相加：
\[
\frac{\partial L}{\partial \bm{W}}
=
\frac{\partial L}{\partial \bm{W}^{(3)}}
+
\frac{\partial L}{\partial \bm{W}^{(5)}}.
\]
前向两次使用同一 $\bm{W}$，更新一步同时改变两处表观层。

\subsection{小结}

参数是带梯度的命名张量；初始化给定 $P$，绑定使若干层共享同一 $\bm{W}$ 并把梯度相加。

\subsection{练习}

在含共享层的感知机上核验 $\bm{W}^{(i)}=\bm{W}^{(j)}$ 是否为同一存储，
以及 $\nabla_{\bm{W}}L$ 是否等于两条路径梯度之和。

\section{延后初始化}

定义网络时可以暂不给出输入维 $d$。
此时第一层权重 $\bm{W}^{(1)}\in\mathbb{R}^{h\times d}$ 中的 $d$ 未知，无法分配。
延后初始化（deferred initialization）把形状推断推迟到数据第一次通过模型。

\subsection{实例化网络}

实例化后、首次前向之前，$\operatorname{params}$ 为空。
设首个批量 $\bm{X}\in\mathbb{R}^{n\times d}$ 到达，由 $d$ 确定
\[
\bm{W}^{(1)}\in\mathbb{R}^{h\times d},
\qquad
\bm{W}^{(2)}\in\mathbb{R}^{q\times h},
\]
再依层顺序分配并抽样。
仅第一层真正依赖未知输入维，但框架仍按拓扑顺序完成全部形状。

\subsection{小结}

首次前向把未知维代入，使
\[
\operatorname{shape}(\bm{W}^{(\ell)})
\]
全部成为已知后再初始化。

\subsection{练习}

核验：只给出第一层输入维时哪些层立即分配；维不匹配是否使乘法
$\bm{W}\bm{x}$ 无定义；输入维改变时共享参数如何约束形状。

\section{自定义层}

层是块的特例。当预置层不包含所需映射时，按同一接口定义新层：
前向 $f$ 加上可选的 $\operatorname{params}(f)$。

\subsection{不带参数的层}

中心化层是无参数映射
\[
f(\bm{x})
=
\bm{x}-\bar{x}\bm{1},
\qquad
\bar{x}=\frac{1}{n}\sum_{i=1}^{n}x_i.
\]
对批量矩阵按约定的轴取均值后广播相减。
浮点下 $\sum_i f(\bm{x})_i$ 只是接近 $0$。
该层可嵌入任意 $\operatorname{Seq}$。

\subsection{带参数的层}

自定义全连接层需要 $\bm{W}\in\mathbb{R}^{q\times d}$、$\bm{b}\in\mathbb{R}^{q}$，
\[
f(\bm{x})
=
\operatorname{ReLU}(\bm{W}\bm{x}+\bm{b}).
\]
用框架的参数构造器登记 $\bm{W},\bm{b}$，则访问、初始化、共享、序列化都沿用统一接口，
而不必为该层单独写读写程序。

\subsection{小结}

自定义层是映射 $f$，可无参数或带 $\bm{W},\bm{b}$；登记后即可出现在任意架构中。

\subsection{练习}

核验二次型层
\[
y_k=\sum_{i,j}W_{ijk}\,x_i x_j
\]
的形状 $k\mapsto y_k$，以及把输入映到傅里叶系数前半部分的线性映射。

\section{读写文件}

训练得到 $\bm{\theta}$ 后，需要把张量或整份参数字典写入存储，以便部署或断点续训。

\subsection{加载和保存张量}

单个张量的读写是一对互逆映射
\[
\operatorname{save}(\bm{X},p),\qquad
\bm{X}'=\operatorname{load}(p),
\]
在格式一致时 $\bm{X}'=\bm{X}$。
列表 $(\bm{X}_1,\ldots,\bm{X}_m)$ 与字典 $k\mapsto\bm{X}_k$ 同样可往返，
后者直接对应“全部权重名到张量”。

\subsection{加载和保存模型参数}

整网保存的是参数字典而非可执行架构，因为架构含任意代码，难以序列化。
恢复过程为
\[
f\leftarrow \operatorname{arch}(),
\qquad
\bm{\theta}(f)\leftarrow \operatorname{load}(p).
\]
必须先用代码重建与保存时相同的 $f$，再填入 $\bm{\theta}$。
随机初始化被文件中的值覆盖。

\subsection{小结}

$\operatorname{save}/\operatorname{load}$ 作用于张量；整网只固化 $\bm{\theta}$，架构仍由代码给出。

\subsection{练习}

核验只加载前两层 $\bm{\theta}_{1:2}$ 并入新架构是否形状相容，
以及同时保存架构时必须对 $f$ 施加何种可序列化限制。

\section{GPU}

张量带有设备（device）。默认在主机内存由 CPU 计算；
GPU 把同一线性代数核放到高吞吐器件上。
记设备集合 $\{\mathrm{cpu},\mathrm{cuda}{:}0,\ldots,\mathrm{cuda}{:}{k-1}\}$。

\subsection{计算设备}

$\mathrm{cpu}$ 表示全部主机核与内存；$\mathrm{cuda}{:}i$ 表示第 $i$ 块卡及其显存（$i$ 从 $0$ 起）。
$\mathrm{cuda}{:}0$ 与 $\mathrm{cuda}$ 等价。
可用卡数 $k$ 查询后，第 $i$ 块记为 $\mathrm{cuda}{:}i$。

\subsection{张量与GPU}

每个张量 $\bm{X}$ 有位置 $\operatorname{loc}(\bm{X})$。
二元运算 $\bm{X}+\bm{Y}$ 有定义当且仅当
\[
\operatorname{loc}(\bm{X})=\operatorname{loc}(\bm{Y}),
\]
否则结果的存储位置与计算位置均不确定。

\subsubsection{存储在GPU上}

创建时可指定 $\operatorname{loc}(\bm{X})=\mathrm{cuda}{:}0$。
该张量只占用该卡显存，容量受显存上界约束。
第二块卡上的张量满足 $\operatorname{loc}(\bm{Y})=\mathrm{cuda}{:}1$。

\subsubsection{复制}

若 $\operatorname{loc}(\bm{X})\neq\operatorname{loc}(\bm{Y})$，应先做设备间复制
\[
\bm{Z}=\operatorname{copy}_{\operatorname{loc}(\bm{Y})}(\bm{X}),
\]
再计算 $\bm{Z}+\bm{Y}$。
若 $\bm{Z}$ 已在目标设备，再次请求同一设备不分配新存储，即 $\operatorname{id}$ 不变。

\subsubsection{旁注}

器件间传输带宽远低于器件内计算。
多次小复制的开销通常大于一次大复制。
打印或转为主机数组会把数据拉回内存，并可能触发全局解释器锁，从而阻塞全部 GPU 核。

\subsection{神经网络与GPU}

模型参数同样带位置。把 $\bm{\theta}$ 放到设备 $d$ 后，输入必须满足
\[
\operatorname{loc}(\bm{X})=d,
\]
前向才在 $d$ 上计算 $f_{\bm{\theta}}(\bm{X})$。
数据和参数同设备是有效训练的成立条件。

\subsection{小结}

运算要求全部输入同设备；默认 $\operatorname{loc}=\mathrm{cpu}$。
不必要的 $\operatorname{copy}$ 会支配 $T$。

\subsection{练习}

比较大矩阵乘法上 CPU 与 GPU 的 $T$，并核验逐次把 Frobenius 范数拉回主机
是否显著增加 $T_{\mathrm{copy}}$；两卡并行与单卡顺序的时间比应接近线性。
